% ============================================================
% main.tex — selve prøven
%
% Denne fila er struktur og innhold. Alle typografiske
% valg bor i provestil.sty, som også laster babel.
%
% Kompiler med pdflatex:
%   latexmk -pdf main
% (to kjøringer: «Side 2 av 5» trenger sidetallet fra
% forrige runde — latexmk ordner det selv.)
% ============================================================


\documentclass[12pt]{article}


% ============================================================
% LAST INN PRØVESTIL
%
% [forside]: tittelblokk og instrukser på en egen side uten
% sidetall; Del 1 begynner på side 1. Til heldagsprøver og
% prøver med to deler. Uten valget står tittelblokka øverst
% på side 1 med oppgavene rett under — til mindre prøver.
%
% [hefte]: som [forside], men lagt til rette for tosidig
% utskrift og falsing — tom bakside på forsiden, hver del på
% en høyreside, og sidetall delelig med 4.
%
% [statussjekk]: tittel, undertittel og dato — ingen felt.
% Del oppgavene i \nivaa[...]{1}, {2}, {3} i stedet for \del.
%
% [innrykk]: avsnittsinnrykk i stedet for luft mellom
% avsnittene. Valgene kan kombineres.
% ============================================================
\usepackage[forside]{provestil}
% \usepackage[hefte]{provestil}
% \usepackage{provestil}
% \usepackage[statussjekk]{provestil}
% \usepackage[innrykk]{provestil}


% ============================================================
% EGENDEFINERTE PAKKER
%
% Legg til prosjektspesifikke pakker her.
% Ikke sett pakker andre steder i filen.
% ============================================================

% --- Matematikk (standard) ---
\usepackage{amsmath, amssymb}

% --- Vektorer: \vv{F}, flat pil over symbolet. Indekser utenfor: \vv{F}_1 ---
\usepackage[e]{esvect}

% --- Enheter: \qty{9.81}{\metre\per\second\squared} ---
\usepackage{siunitx}
% Norsk oppsett (desimalkomma, prikk i standardform, «til», «og»)
% ligger i stilfila. Overstyr her ved behov:
% \sisetup{exponent-product=\times}

% --- Figurer ---
\usepackage{graphicx}
\graphicspath{{figurer/}}

% --- Annet ---
\usepackage{tikz}


% ============================================================
% EGENDEFINERTE KOMMANDOER
% ============================================================

% Oppreist d:
% \newcommand{\dd}{\mathrm{d}}
% \newcommand{\dv}[2]{\frac{\dd #1}{\dd #2}}


% ============================================================
% METADATA
%
% Tomme felt vises som tomme (ikke utelatt). \navnetikett
% bytter «Navn» mot f.eks. «Kandidatnummer».
% ============================================================
\title{Prøve i Fysikk 1}
\undertittel{Kapittel 3 og 4: Krefter og bevegelse}
\date{Torsdag 15. oktober 2026}
\tid{90 minutter}
\hjelpemidler{Del 1: ingen. Del 2: alle, unntatt kommunikasjon}
% \navnetikett{Kandidatnummer}


% ============================================================
% DOKUMENT
% ============================================================
\begin{document}

\maketitle

% Med [forside] står alt herfra til første \del på forsiden.
Svar på alle oppgavene. Vis utregning og begrunn svarene;
et riktig svar uten begrunnelse gir ikke full uttelling.
Bruk gjerne figurer.

Del 1 leveres inn før du begynner på Del 2. Prøven har
5 oppgaver og gir til sammen 20 poeng.


% ============================================================
% DEL 1
%
% Med [forside] lukker første \del forsiden og begynner på
% side 1. Sløyf \del-linjene helt på en mindre prøve uten
% deler (og uten [forside]) — oppgavene telles likt uansett.
% ============================================================
\del[Uten hjelpemidler]{1}


\begin{oppgave}[Krefter på skråplan \poeng{4}]
En kloss med masse $m$ ligger i ro på et skråplan som danner
vinkelen $\theta$ med horisontalen.

\begin{deloppgaver}
  \item Tegn figur og sett på alle kreftene som virker på klossen.
  \item Vis at normalkraften er $N = mg\cos\theta$.
  \item Hvor stor må friksjonskraften minst være for at klossen
        skal ligge i ro?
\end{deloppgaver}
\end{oppgave}


\begin{oppgave}[Bevegelseslikningene \poeng{4}]
En ball kastes rett opp med farten $v_0 = \qty{12}{\metre\per\second}$.
Se bort fra luftmotstand, og bruk $g = \qty{9.8}{\metre\per\second\squared}$.

\begin{deloppgaver}
  \item Hvor høyt kommer ballen?
  \item Hvor lang tid tar det før ballen er tilbake i utgangspunktet?
\end{deloppgaver}
\end{oppgave}


\begin{oppgave}[Vektorer \poeng{2}]
To krefter $\vv{F}_1$ og $\vv{F}_2$ virker på samme punkt.
$\vv{F}_1$ har størrelse \qty{3.0}{\newton} og peker mot øst;
$\vv{F}_2$ har størrelse \qty{4.0}{\newton} og peker mot nord.
Hvor stor er den resulterende kraften?
\begin{flervalg}
  \item \qty{1.0}{\newton}
  \item \qty{5.0}{\newton}
  \item \qty{7.0}{\newton}
  \item \qty{12}{\newton}
\end{flervalg}
\end{oppgave}


% ============================================================
% DEL 2
%
% \clearpage om Del 1 skal leveres inn før Del 2 begynner.
% Skal oppgavene starte på 1 igjen her:
% \setcounter{oppgaveinner}{0} etter \del.
% ============================================================
\clearpage
\del[Med hjelpemidler]{2}


\begin{oppgave}[Bilen som bremser \poeng{6}]
En bil kjører med farten \qty{80}{\kilo\metre\per\hour} når
sjåføren oppdager et hinder \qty{50}{\metre} lenger framme.
Reaksjonstiden er \qty{0.8}{\second}, og bremsene gir en
konstant retardasjon på \qty{6.5}{\metre\per\second\squared}.

\begin{deloppgaver}
  \item Hvor langt kjører bilen i løpet av reaksjonstiden?
  \item Rekker bilen å stoppe før hinderet? Vis utregningen.
\end{deloppgaver}

Figuren under viser et fart--tid-diagram for bevegelsen.

\begin{figure}[ht]
  \centering
  \begin{tikzpicture}[>=stealth, x=1.1cm, y=0.9cm]
    \draw[->,thick] (0,0) -- (5.2,0) node[right,font=\small] {$t$ (s)};
    \draw[->,thick] (0,0) -- (0,3.4) node[above,font=\small] {$v$ (m/s)};
    \draw[gray,thin,dashed] (0.8,0) -- (0.8,2.5);
    \node[below,font=\small,gray] at (0.8,0) {$0{,}8$};
    \node[below,font=\small,gray] at (4.2,0) {$t_{\text{stopp}}$};
    \node[left,font=\small] at (0,2.5) {$22{,}2$};
    \fill[darkorange!15] (0,0) -- (0,2.5) -- (0.8,2.5) -- (4.2,0) -- cycle;
    \draw[darkolive,thick] (0,2.5) -- (0.8,2.5) -- (4.2,0);
  \end{tikzpicture}
  \caption{Fart som funksjon av tid. Arealet under grafen er
           strekningen.}
  \label{fig:vt}
\end{figure}

\begin{deloppgaver}[resume]
  \item Bruk figur~\ref{fig:vt} til å forklare hvordan
        stopplengden kan leses ut av diagrammet.
\end{deloppgaver}
\end{oppgave}


\begin{oppgave}[Numerisk løsning \poeng{4}]
Programmet under simulerer fallet til en kule med luftmotstand
ved hjelp av Eulers metode.

\begin{python}
import numpy as np

m, g, k = 0.10, 9.81, 0.02     # masse, g, luftmotstand
v, t, dt = 0.0, 0.0, 0.01

while t < 5.0:
    a = g - (k / m) * v**2      # akselerasjon (m/s^2)
    v = v + a * dt
    t = t + dt

print(f"Fart etter {t:.1f} s: {v:.2f} m/s")
\end{python}

\begin{deloppgaver}
  \item Forklar hva linje 7 gjør, og hvilken fysisk lov den bygger på.
  \item Hva vil programmet skrive ut hvis $k$ settes til null?
        Begrunn svaret uten å kjøre programmet.
  \item Endre programmet slik at det også regner ut hvor langt
        kula har falt.
\end{deloppgaver}
\end{oppgave}

\separator

\begin{center}
\sffamily Lykke til!
\end{center}

\end{document}
